% Subseção da F-ROTA (trajeto mais rápido), com as extensões F-ROTA-2 (A* e ponto no mapa) e
% F-ROTA-OPC (Yen) para caber no orçamento de páginas. Dono: Henrique.
% Fontes: docs/plano-de-desenvolvimento/features/feature-F-ROTA.md, feature-F-ROTA-2.md e
% feature-F-ROTA-OPC.md (Como funciona e Resultados); números de
% benchmarks/medir_modelo_rota.py e benchmarks/medir_rota_extensoes.py em 18/11/2026.
% Figura: mapa da página Trajetos (mapa_trajeto) para dois pontos escolhidos por coordenadas.
\subsection{Trajeto mais rápido}\label{sec:F-ROTA}

Dadas duas paradas e uma faixa horária, a ferramenta acha o trajeto que chega mais cedo em
média, somando caminhada, espera e viagem, e o mostra no mapa por trecho: a pé, espera e
linha usada, com os tempos (Figura~\ref{fig:trajetos-mapa}). A visão é $G_E$: como trocar de
linha obriga a desembarcar e pagar um novo embarque, que pesa a espera $H/2$ da nova linha
(F-ESPERA), a baldeação vira custo comum do caminho e basta um Dijkstra com o heap próprio,
sem penalidade arbitrária. A busca para quando o destino sai do heap, o caminho é
reconstruído por arestas, porque o tipo de cada passo está na aresta, e as arestas são
agrupadas em trechos. O custo é $O((V + E)\log V)$ por consulta.

\begin{figure}[t]
\centering
\includegraphics[width=0.46\linewidth]{trajetos-mapa}
\caption{Trajeto do centro (Av.\ Afonso Pena) a Venda Nova no $G_E$ do pico da manhã,
escolhido por coordenadas: linha 4103, troca a pé (círculo laranja) e linha 63; 50~min~09~s
entre as paradas e 51~min~29~s de porta a porta, com 1 baldeação.}
\label{fig:trajetos-mapa}
\end{figure}

No feed real (18/11/2026, 9.660 paradas, 698 padrões), o $G_E$ do pico da manhã tem
$V = 46.102$ e $E = 209.025$. Em 200 pares de paradas sorteados, a consulta levou
$65{,}9$~ms em média e $129{,}6$~ms no pior caso (Intel Core i5-13600KF), dentro do limite
de 1~s do RNF-05; todos tiveram trajeto, com tempo mediano de 61,1~min e 0 (30 pares),
1 (89), 2 (65) e 3 ou mais (16) baldeações, e o custo coincidiu com o do NetworkX nos 200.

\emph{Ponto no mapa e A*.} Um clique ou um par de coordenadas é ligado à parada mais próxima
por varredura linear, $O(V)$, em 4,90~ms por ponto, e a caminhada até ela entra no tempo
porta a porta. A página também calcula o trajeto com A*, usando como estimativa a distância
em linha reta até o destino dividida pela maior velocidade entre as arestas de $G_E$. Essa
estimativa é consistente, $h(u) \le w(u,v) + h(v)$ pela desigualdade triangular, e o A*
devolveu o mesmo custo do Dijkstra nos 200 pares. O ganho, porém, foi quase nulo: expandiu
99,7\% dos vértices do Dijkstra (mediana) e levou 90,8~ms contra 72,3~ms, porque paga uma
distância por vértice. Dois motivos: 181 trechos com o mesmo horário em duas paradas
seguidas no GTFS recebem o piso de 1~s e levam a velocidade máxima a 1.258~m/s, o que deixa
a estimativa muito pequena; e ela não enxerga as esperas, que são boa parte do custo. Usar
uma velocidade realista daria mais poda, mas tiraria a garantia de ótimo; o grupo manteve o
A* admissível e registra a comparação como resultado.

\emph{Trajetos alternativos.} O algoritmo de Yen lista os $k$ trajetos sem ciclo de menor
tempo: cada vértice do trajeto anterior vira ponto de desvio, as saídas já usadas a partir
da mesma raiz são desativadas pela API de falhas do núcleo, e um A* completa o desvio. O
custo é $O(k\,n\,(V + E)\log V)$, com $n$ vértices no maior trajeto, e o tempo cresce com
$k$: mediana de 2,3~s com $k = 3$ e 4,6~s com $k = 5$ em 20 pares, até 25,4~s, por isso a
página só calcula sob demanda. O resultado coincidiu com o do NetworkX em redes aleatórias.

\emph{Limitações:} o tempo é o esperado para quem chega à parada num instante qualquer da
faixa (espera $H/2$), não o de um horário de saída exato; a caminhada é em linha reta, e o
ponto escolhido no mapa só é ligado à parada mais próxima. Alternativos distintos diferem em
pelo menos uma aresta de $G_E$ e podem ser variações pequenas do mesmo trajeto.
